求解得到全天计划购电量 $57526.2435\,\mathrm{kWh}$，购电费 $33801.4955$ 元。无储能基准购电费为 $48052.0466$ 元，储能使成本减少 $14250.5510$ 元，降幅为 $29.6565\%$。总充电量和总放电量分别为 $19842.7105\,\mathrm{kWh}$ 和 $17858.4394\,\mathrm{kWh}$，弃光量为 0。指定时段购电结果列于表\ref{tab:q1-purchase}，六个四小时区间的充放电结果及首末储电量列于表\ref{tab:q1-storage}，表头和排列沿用题目表 1、表 2；电量单位为 kWh，购电费单位为元。
\begin{table}[htbp]
\centering
\caption{微网在指定时间段的购电量及全天的购电量和购电费}\label{tab:q1-purchase}
\normalsize\setlength{\tabcolsep}{5pt}
\begin{tabular}{@{}cccccc@{}}
\toprule
时间段 & 购电量 & 时间段 & 购电量 & 时间段 & 购电量 \\
\midrule
10:00-10:10 & 0.0000 & 12:00-12:10 & 486.4029 & 14:00-14:10 & 0.0000 \\
16:00-16:10 & 394.9315 & 18:00-18:10 & 636.9826 & 20:00-20:10 & 0.0000 \\
全天购电量 & \multicolumn{2}{c}{57526.2435} & 全天购电费 & \multicolumn{2}{c}{33801.4955} \\
\bottomrule
\end{tabular}
\end{table}

\begin{table}[htbp]
\centering
\caption{储能设备在指定时间段的充放电量及0:00和24:00的储电量}\label{tab:q1-storage}
\normalsize\setlength{\tabcolsep}{5pt}
\begin{tabular}{@{}cccccc@{}}
\toprule
时间段 & 充电量 & 放电量 & 时间段 & 充电量 & 放电量 \\
\midrule
0:00-4:00 & 4226.3109 & 0.0000 & 4:00-8:00 & 833.3333 & 6956.9814 \\
8:00-12:00 & 4448.6339 & 1702.9970 & 12:00-16:00 & 5274.7881 & 91.1014 \\
16:00-20:00 & 0.0000 & 6247.4915 & 20:00-24:00 & 5059.6443 & 2859.8681 \\
0:00 储电量 & \multicolumn{2}{c}{6000.0000} & 24:00 储电量 & \multicolumn{2}{c}{6000.0000} \\
\bottomrule
\end{tabular}
\end{table}


图\ref{fig:q1-dispatch}显示，储能主要在低价时段和光伏出力较高的时段充电，在电价上升后释放能量；中午光伏提高时，外网购电明显下降。储电量始终位于 $1200\sim10800\,\mathrm{kWh}$，并于日末回到 $6000\,\mathrm{kWh}$，说明成本下降来自时段间能量转移，而非透支初始储能。最大购电功率为 $8458.8273\,\mathrm{kW}$；该值依赖“外网接口容量未给定、暂不设上限”的建模边界。
\begin{figure}[htbp]
\centering
\includegraphics[width=0.94\linewidth]{figures/q1/q1_figure_2.png}
\caption{问题一最优调度与储能状态}\label{fig:q1-dispatch}
\end{figure}
\FloatBarrier
